% ══════════════════════════════════════════════════════════════
%  SECTION 7: CONCLUSION
% ══════════════════════════════════════════════════════════════

\section{Conclusion}
\label{sec:conclusion}

FIFA's 2022 stoppage-time directive provides a rare natural experiment
in which the effective duration of a competitive interaction changed
while rules, incentives, and participants remained constant. Three
findings emerge.

First, the directive produced a dose-response gradient in enforcement
intensity that tracks institutional proximity to FIFA with remarkable
precision. The gradient is itself the first-stage result: it documents
how a global regulator's signal propagates through a decentralised
governance architecture, amplified by professional referee networks and
broadcast media, dampened by institutional insularity.

Second, the additional minutes produced goal displacement, not goal
creation. Total goals were unchanged; goals migrated from the final
minutes of regulation into expanded stoppage time, raising the share
of matches with a decisive stoppage-time moment from 7.3\% to 10.3\%.
A simple model of optimal attacking intensity under a fatigue
constraint rationalises this pattern: when the trailing team's time
horizon expands, it substitutes patience for desperation, spreading
effort without increasing aggregate output. The match exhibited
approximately constant returns to time at the margin studied.

Third, compliance was achieved through focal-point coordination rather
than monitoring or sanctions. The World Cup demonstration crystallised a
new behavioural norm that became self-reinforcing: 95\% of referees
individually increased their stoppage time, with no decay over three
seasons. The pre-directive variance in stoppage time reflected a
coordination failure, not genuine preference heterogeneity.

Several extensions await future work. The model's prediction that
displacement should be stronger in close matches is testable with
finer data on scoring rates by game state and time bin. The directive's
implications for player welfare---particularly whether clubs' squad
rotation response (fewer matches per season despite longer individual
matches) was optimal or imposed costs on marginal squad
members---merit sustained attention as longer matches become the norm.
The betting-market response, which appears rapid and well-calibrated in
preliminary analysis, deserves formal treatment as a test of market
efficiency under regime change. Finally, the dose-response gradient
invites structural modelling of how enforcement norms diffuse through
hierarchical sporting governance, with potential applications to other
settings where global regulators issue guidance that national bodies
implement with discretion.

